
\begin{abstract}
 We study Boolean Holant problems with complex-valued signature sets closed under conjugation. Such sets arise naturally in tensor-network expressions for classical strong simulation of quantum circuits. We prove a complexity dichotomy for such problems with an explicit tractability criterion. This extends the dichotomy for real-valued Holant problems, with the same four tractability conditions.


Our proofs use Xia's projective binary group framework and quantum entanglement theory.
The conjugate closure assumption precisely makes $k$-uniformity, directly applicable
to the classification of Holant problems, by realizing reduced density matrices via Holant gadgets. 
We also use the classification of absolutely maximally entangled states to resolve a particular $6$-ary obstruction in our
inductive proof of the \#P-hardness. 

\end{abstract}

\begin{comment}

We prove a complexity dichotomy for complex Boolean Holant problems defined by signature sets which are closed under function-wise conjugation.
This extends the dichotomy for real Boolean Holant problems by Shao and Cai on FOCS 2020 \cite{realholant}.
It is proved that $\Holant(\mathcal{F})$ is either tractable or \#P-hard, for an arbitrary complex signature set $\mathcal{F}$ such that $\mathcal{F}=\overline{\mathcal{F}}$.\footnote{We assume that every signature $f$ in $\mathcal{F}$ takes algebraic complex value.}
The dichotomy has an explicit criteria.
The proof utilizes the framework of Xia \cite{MingjiProgram}, who suggests classifying all problems $\Holant(\mathcal{F})$ with unknown complexity according to the projective binary group generated by $\mathcal{F}.$
We also use results from representation theory and quantum information theory, specifically, the irreducible representations of polyhedral groups and the classification of absolutely maximally entangled states.
    
\end{comment}
